\section{Aggregate Lower Bound Theorem}
\label{sec:lower_bound}

\subsection{Event-local decomposition}

\begin{definition}[Event-Local Bundle Decomposition]
\label{def:bundle_decomp}
For each revealed-sacrifice event $(i_n, X_n, Y_n, t_n)$, define
the \emph{event-local} decomposition coefficients
$\alpha_{n,c} \geq 0$ for each capability $c \in Y_n$,
satisfying:
\[
  \sum_{c \in Y_n} \alpha_{n,c} = 1
  \qquad\text{(partition of unity within the bundle)}.
\]
The coefficients represent the share of bundle $Y_n$ attributable
to capability~$c$, determined at the time of event~$n$ and
\emph{fixed thereafter}.  Each event carries its own
decomposition, independent of other events.  The decomposition can
be obtained from prior hedonic
regression~\citep{rosen1974hedonic} on similar trades, domain
knowledge, or the equal-weight default
$\alpha_{n,c} = 1/|Y_n|$.
\end{definition}

\begin{remark}[Event-local versus global regression]
Classical hedonic regression~\citep{rosen1974hedonic} fits a
\emph{global} model across all observed trades simultaneously,
producing the tightest decomposition but breaking the monotonicity
of the max-attribution bound
(Proposition~\ref{prop:monotone}) when new observations refit all
coefficients.  The event-local formulation fixes coefficients at
observation time, sacrificing fit quality for the monotonicity
guarantee.
\end{remark}

\subsection{Central result}

\begin{theorem}[Aggregate B-to-C Lower Bound]
\label{thm:aggregate_bound}
For a sequence of revealed-sacrifice events
$\{(i_n, X_n, Y_n, t_n)\}_{n=1}^N$ with bundles jointly covering
a subset $U$ of unbenchmarked capability-space, under
Assumptions~S0--S5 of
Section~\ref{sec:assumptions}, the theorem has two levels:

\medskip\noindent
\textbf{Part~A (Bundle-level lower bound, requires S0--S2 only).}
For each event~$n$:
\[
  \volR(Y_n) \;\geq\; \Delta\volL(X_n).
\]
For $N$ events with \emph{pairwise poset-independent} bundles:
\[
  \volR^U \;\geq\; \sum_{n=1}^{N} \Delta\volL(X_n),
\]
where $\volR^U$ is the realized capability volume restricted to
the unbenchmarked subset $U$ covered by the union of bundles.

\medskip\noindent
\textbf{Part~B (Per-capability refinement, requires S0--S5).}
When bundles overlap, define for each capability $c \in U$:
\[
  \volRlower(c) \;=\; \max_{n:\, c \in Y_n}
  \alpha_{n,c} \cdot \Delta\volL(X_n).
\]
Then, provided the capabilities $c \in U$ are pairwise
poset-independent:
\[
  \volR^U \;\geq\; \sum_{c \in U} \volRlower(c).
\]
\end{theorem}

\begin{proof}[Proof sketch]
\emph{Part~A.}
Step~1: For each event~$n$, the calibration chain
(S0 + S1 + S2) gives
$\volR(Y_n) \geq \Ui[_{n}](Y_n) \geq \Ui[_{n}](X_n) \geq
\Delta\volL(X_n)$.
Step~2: For poset-independent bundles, the sub-posets $\{Y_n\}$
are disconnected components.  By axiom~M4 (additivity on
poset-disjoint
components,~\citep[Proposition~1]{lasser2026poset}),
$\volR(\bigcup Y_n) = \sum_n \volR(Y_n) \geq \sum_n
\Delta\volL(X_n)$.

\emph{Part~B.}
Step~3: S4(b) decomposes
$\volR(Y_n) = \sum_c \volR(c)$.
S5 gives
$\alpha_{n,c} \leq \volR(c)/\volR(Y_n)$,
hence $\volR(c) \geq \alpha_{n,c} \cdot \volR(Y_n) \geq
\alpha_{n,c} \cdot \Delta\volL(X_n)$.
Step~4: Max-attribution across overlapping events.
Step~5: Global summation over poset-independent capabilities
via~M4.
\end{proof}

\begin{remark}[Why poset-independence, not set-disjointness]
\label{rem:poset_indep}
Set-disjointness ($Y_1 \cap Y_2 = \emptyset$) ensures no
\emph{capability} appears in two bundles, but does not rule out
subsumption or cooperative links across bundles.  If a cooperative
capability $c_{\mathrm{coop}}$ depends on capabilities in both
$Y_1$ and $Y_2$, its $\volR$ contribution is not attributable to
either bundle alone.  Poset-independence rules this out: under
poset-independence, the sub-posets are disconnected components
and $\volL$ (hence $\volR$) decomposes additively by
axiom~M4~\citep[Proposition~1]{lasser2026poset}.
\end{remark}

\subsection{B-to-C ratio and monotonicity}

\begin{definition}[B-to-C Ratio Under Revealed Sacrifice]
\label{def:btoc_ratio}
The B-to-C ratio under revealed-sacrifice observation is:
\[
  \betaL(G, T) \;=\; \volRlower(G, T) \;/\; \volL(G),
\]
where $\volRlower$ is the aggregate lower bound from
Theorem~\ref{thm:aggregate_bound} and $T$ indexes the trade
window.  $\betaL \in [0, 1]$ because $\volR \leq \volL$ by
definition (the exercised sub-poset is a subset of the full
poset).
\end{definition}

\begin{proposition}[Monotone Accumulation Under Max-Attribution]
\label{prop:monotone}
Define the per-capability lower bound as the maximum attribution
across all observed events:
\[
  \volRlower(c) \;=\; \max_{n:\, c \in Y_n}
  \alpha_{n,c} \cdot \Delta\volL(X_n),
\]
and the aggregate as $\volRlower = \sum_{c \in U} \volRlower(c)$.
Then the aggregate is non-decreasing in the number of observed
events: for $N' > N$,
\[
  \volRlower(N') \;\geq\; \volRlower(N),
\]
provided the additional events satisfy the applicable assumptions.
\end{proposition}

\begin{proof}[Proof sketch]
Each per-capability bound is a maximum over a growing set of
candidates.  Adding events can only weakly increase a max.  The
aggregate is a sum of non-decreasing terms, hence non-decreasing.
\end{proof}
